\subsection{Preparazione del dataset}

Le partite utilizzate per l'addestramento, disponibili in formato SGF, vengono elaborate per ottenere un dataset di esempi supervisionati. Ogni esempio è costituito dalla rappresentazione dello stato della scacchiera in un determinato momento della partita e dalla mossa successivamente effettuata dal giocatore di turno.

La pipeline è implementata principalmente nei moduli \texttt{go\_preprocessing.py} e \texttt{go\_dataset.py} e comprende il parsing delle partite, l'estrazione delle feature e delle mosse target, la suddivisione dei dati e il loro salvataggio.

\subsubsection{Lettura delle partite ed estrazione degli esempi}

La lettura dei file SGF viene effettuata mediante la libreria \texttt{sgfmill}, utilizzata per ricostruire la sequenza delle mosse. Durante il parsing viene verificata la dimensione della scacchiera, così da escludere le partite non compatibili con la configurazione in elaborazione.

Per ciascuna partita, il programma ricostruisce progressivamente lo stato del goban. Prima di applicare ogni mossa vengono memorizzati la configurazione corrente della scacchiera e il colore del giocatore di turno: queste informazioni costituiscono rispettivamente l'input e l'etichetta dell'esempio.

La mossa target viene codificata mediante un indice intero. Le intersezioni assumono valori compresi tra \(0\) e \(N^2-1\), mentre l'azione \textit{PASS} è rappresentata dall'indice \(N^2\), dove \(N\) indica la dimensione della scacchiera.

Durante la ricostruzione vengono controllate anche le mosse illegali. Gli esempi corrispondenti a mosse non applicabili vengono esclusi e la partita viene segnalata come non valida.

Una gestione specifica riguarda la conclusione delle partite: due mosse \textit{PASS} consecutive determinano la terminazione del gioco, ma alcuni file SGF possono non riportarle entrambe esplicitamente. Quando il risultato registrato è compatibile con una conclusione basata sul punteggio, il preprocessing ricostruisce gli eventuali PASS mancanti, evitando di introdurre mosse artificiali nei casi di abbandono o conclusione per tempo.

\subsubsection{Estrazione delle feature}

La funzione \texttt{extract\_features()} converte lo stato della partita in un tensore numerico di forma \(7\times N\times N\), coerente con l'input richiesto dal modello.

La rappresentazione considera tre configurazioni temporali della scacchiera, quella corrente (\(T\)) e le due precedenti (\(T-1\) e \(T-2\)). Per ciascuna vengono utilizzati due canali che distinguono le pietre del giocatore corrente da quelle dell'avversario; un ulteriore canale identifica il colore del giocatore di turno.

I primi sei canali sono mappe binarie, mentre il settimo assume valore costante su tutta la scacchiera: 1 per il Nero e 0 per il Bianco.

La costruzione dei canali avviene rispetto al giocatore corrente, mantenendo una prospettiva uniforme indipendentemente dal colore che deve effettuare la mossa. Il principio di implementazione è il seguente:

\begin{verbatim}
player_channel = history_idx * 2
opponent_channel = history_idx * 2 + 1

tensor[player_channel] = board_array == current_color
tensor[opponent_channel] = board_array == opposite_color(current_color)
\end{verbatim}

Il tensore risultante viene convertito in formato \texttt{float32}, mentre la mossa target viene memorizzata separatamente come indice intero. Ogni esempio presenta pertanto un input di forma \([7,N,N]\) e un'etichetta scalare.

\subsubsection{Suddivisione del dataset}

Gli esempi vengono organizzati in tre insiemi: training, validation e test, secondo una proporzione rispettivamente pari all'80\%, 10\% e 10\%.

La suddivisione avviene a livello di partita e non di singola posizione. Questa scelta evita che esempi appartenenti alla stessa partita siano distribuiti tra insiemi differenti, riducendo il rischio di data leakage e rendendo più significativa la valutazione della capacità di generalizzazione.

La selezione viene effettuata in maniera riproducibile mediante un generatore casuale inizializzato con seed pari a 42.

Per il dataset 19$\times$19, considerate le dimensioni dell'archivio originale, è inoltre possibile limitare il numero di partite elaborate. Nella configurazione utilizzata sono state selezionate al massimo 10.000 partite, prima della suddivisione nei tre insiemi. Anche questa selezione è deterministica.

\subsubsection{Data augmentation mediante simmetrie \texorpdfstring{\(D_4\)}{D4}}

Durante l'addestramento viene applicata una tecnica di \textit{data augmentation} basata sulle otto simmetrie del gruppo diedrale \(D_4\), corrispondenti alle quattro rotazioni di \(0^\circ\), \(90^\circ\), \(180^\circ\) e \(270^\circ\) e alle rispettive riflessioni.

Queste trasformazioni mantengono inalterate le relazioni tra le intersezioni e permettono di ottenere configurazioni equivalenti del gioco con orientamenti differenti.

La trasformazione viene applicata sia alle feature sia alla mossa target, che deve essere rimappata nella nuova posizione. L'azione \textit{PASS}, non essendo associata a un'intersezione, rimane invece invariata.

L'implementazione è contenuta nella funzione \texttt{apply\_board\_symmetry()} del modulo \texttt{go\_dataset.py}. A ogni campione viene applicata casualmente una delle otto trasformazioni durante il caricamento, senza generare preventivamente copie aggiuntive del dataset.

La data augmentation è attiva esclusivamente per il training; validation e test mantengono la rappresentazione originale, garantendo la riproducibilità delle metriche.

\subsubsection{Salvataggio e caricamento dei dati}

I dataset vengono salvati in formato NumPy, attraverso il formato compresso \texttt{NPZ} oppure mediante file \texttt{NPY}, utilizzati nelle configurazioni principali per gestire gli insiemi di dimensioni maggiori.

In quest'ultimo caso, ciascun insieme contiene due file:

\begin{itemize}
    \item \texttt{X.npy}, contenente i tensori delle feature in formato \texttt{float32};
    \item \texttt{y.npy}, contenente gli indici delle mosse target in formato \texttt{int64}.
\end{itemize}

Per evitare di caricare interamente in memoria dataset particolarmente grandi, viene utilizzato il \textit{memory mapping} di NumPy:

\begin{verbatim}
self.X = np.load("X.npy", mmap_mode="r")
self.y = np.load("y.npy", mmap_mode="r")
\end{verbatim}

Questa modalità consente di accedere ai dati direttamente dai file, caricando in memoria soltanto le porzioni necessarie. Il vantaggio è particolarmente rilevante per il dataset 19$\times$19.

Durante l'addestramento, la classe \texttt{GoDataset} gestisce l'accesso ai campioni, mentre i \textit{DataLoader} di PyTorch organizzano il caricamento in batch. Le trasformazioni casuali vengono applicate esclusivamente agli esempi di training.

La pipeline produce quindi tre insiemi distinti, pronti per l'addestramento e la valutazione, mantenendo la stessa struttura logica per entrambe le dimensioni di scacchiera.
